\documentclass[12pt,a4paper]{amsart}

\usepackage[ngerman]{babel}
\usepackage{amsmath,amssymb,amsthm}
\usepackage{mathtools}
\usepackage[colorlinks=true,linkcolor=blue,citecolor=blue]{hyperref}
\usepackage{geometry}
\geometry{margin=2.5cm}

% -------------------------------------------------------
% Theorem-Umgebungen
% -------------------------------------------------------
\newtheorem{theorem}{Satz}[section]
\newtheorem{lemma}[theorem]{Lemma}
\newtheorem{korollar}[theorem]{Korollar}
\newtheorem{proposition}[theorem]{Proposition}
\theoremstyle{definition}
\newtheorem{definition}[theorem]{Definition}
\newtheorem{bemerkung}[theorem]{Bemerkung}
\newtheorem{beispiel}[theorem]{Beispiel}

% -------------------------------------------------------
% Eigene Makros
% -------------------------------------------------------
\newcommand{\N}{\mathbb{N}}
\newcommand{\Z}{\mathbb{Z}}
\newcommand{\Q}{\mathbb{Q}}
\newcommand{\Fp}{\mathbb{F}_p}
\DeclareMathOperator{\ord}{ord}
\DeclareMathOperator{\ggT}{ggT}

% -------------------------------------------------------
\begin{document}
% -------------------------------------------------------

\title{Der Wilson-Quotient und Wilson-Primzahlen}

\author{Michael Fuhrmann}
\address{Specialist-Mathematics-Projekt, Build 64}
\email{specialist-maths@localhost}
\date{\today}

\begin{abstract}
Der Satz von Wilson garantiert, dass $(p-1)! + 1$ durch jede Primzahl $p$
teilbar ist. Dies motiviert den \emph{Wilson-Quotienten}
$W_p = \bigl((p-1)!+1\bigr)/p \in \Z$.
Eine Primzahl $p$ heißt \emph{Wilson-Primzahl}, wenn $p^2 \mid (p-1)!+1$ gilt,
äquivalent wenn $p \mid W_p$. Wir beweisen vier Hauptresultate:
(i) eine explizite Kongruenz, die $W_p$ durch die harmonische Summe
$\sum_{j=1}^{p-1} j^{-1} \pmod{p}$ ausdrückt,
(ii) die Äquivalenz der Wilson-Primzahl-Bedingung mit $p^2 \mid (p-1)!+1$,
(iii) einen Zusammenhang mit Bernoulli-Zahlen via Kummers Kongruenz
$W_p \equiv -B_{p-1} \pmod{p}$,
und (iv) ein heuristisches Argument für die Seltenheit von Wilson-Primzahlen.
Die einzigen bekannten Wilson-Primzahlen sind $5$, $13$ und $563$;
keine weiteren Wilson-Primzahlen existieren unterhalb von $2 \times 10^{13}$.
\end{abstract}

\maketitle
\tableofcontents

% ================================================================
\section{Einleitung}
% ================================================================

Der Satz von Wilson besagt $(p-1)! \equiv -1 \pmod{p}$ für jede Primzahl $p$.
Das bedeutet $p \mid (p-1)! + 1$, daher ist die ganze Zahl
\[
  W_p = \frac{(p-1)! + 1}{p}
\]
wohldefiniert. Die Frage, ob $p$ auch $W_p$ teilt, ist der Ausgangspunkt
der Theorie der Wilson-Primzahlen.

Wilson-Primzahlen sind analog zu Wieferich-Primzahlen beim kleinen Satz von
Fermat: Genau wie Wieferich-Primzahlen $a^{p-1} \equiv 1 \pmod{p^2}$ erfüllen
(die „quadrierte" Version von Fermats kleinem Satz), erfüllen Wilson-Primzahlen
die „quadrierte" Version von Wilson. Beide sind extrem selten, und in beiden
Fällen sind nur endlich viele bekannt.

Im gesamten Artikel bezeichnet $p$ eine Primzahl, und alle Rechnungen in
Restklassen gelten modulo der nach $\pmod{}$ geschriebenen Primzahl.

% ================================================================
\section{Der Wilson-Quotient und seine grundlegenden Eigenschaften}
% ================================================================

\begin{definition}[Wilson-Quotient und Wilson-Primzahlen]\label{def:wilson_quotient}
  Für eine Primzahl $p$ ist der \emph{Wilson-Quotient}
  \[
    W_p \;=\; \frac{(p-1)!\,+\,1}{p} \;\in\; \Z.
  \]
  Die Primzahl $p$ heißt \emph{Wilson-Primzahl}, wenn $p \mid W_p$ gilt,
  äquivalent wenn $p^2 \mid (p-1)! + 1$.
\end{definition}

\begin{proposition}[Ganzzahligkeit und erste Werte]\label{prop:grundlegendes}
  $W_p$ ist eine positive ganze Zahl. Die ersten Werte sind:
  \[
    W_2 = 1,\quad W_3 = 1,\quad W_5 = 5,\quad W_7 = 103,\quad W_{11} = 329891.
  \]
\end{proposition}

\begin{proof}
  Nach dem Satz von Wilson gilt $p \mid (p-1)! + 1$, also $W_p \in \Z$.
  Positivität folgt aus $(p-1)! + 1 \geq 1$.
  Die numerischen Werte ergeben sich durch direkte Berechnung:
  \begin{align*}
    W_2 &= \tfrac{1!+1}{2} = 1,\\
    W_3 &= \tfrac{2!+1}{3} = 1,\\
    W_5 &= \tfrac{4!+1}{5} = \tfrac{25}{5} = 5,\\
    W_7 &= \tfrac{6!+1}{7} = \tfrac{721}{7} = 103,\\
    W_{11} &= \tfrac{10!+1}{11} = \tfrac{3628801}{11} = 329891.\qedhere
  \end{align*}
\end{proof}

\begin{theorem}[Wilson-Primzahl-Kriterium]\label{satz:wilson_primzahl_kriterium}
  Eine Primzahl $p$ ist genau dann eine Wilson-Primzahl, wenn
  $W_p \equiv 0 \pmod{p}$, d.\,h.\ $p^2 \mid (p-1)! + 1$.
\end{theorem}

\begin{proof}
  Nach Definition gilt $W_p = \bigl((p-1)!+1\bigr)/p$. Dann gilt:
  \[
    p \mid W_p
    \iff p \mid \frac{(p-1)!+1}{p}
    \iff p^2 \mid (p-1)! + 1. \qedhere
  \]
\end{proof}

% ================================================================
\section{Darstellung durch die halbe harmonische Summe}
% ================================================================

Wir leiten eine explizite Formel für $W_p \pmod p$ her.
Sei $m = (p-1)/2$.  Wir definieren:
\[
  \mu_p \;=\; \frac{(-1)^m (m!)^2 + 1}{p} \;\in\; \Z,
  \qquad
  S_p \;=\; \sum_{j=1}^{m} j^{-1} \pmod p,
\]
wobei $j^{-1}$ das multiplikative Inverse von $j$ modulo $p$ bezeichnet.

\begin{bemerkung}
  Die \emph{volle} harmonische Summe $\sum_{j=1}^{p-1} j^{-1} \equiv 0 \pmod{p}$
  gilt stets, da $\{j^{-1} \colon 1 \leq j \leq p-1\} = \{1, \ldots, p-1\}$
  als Menge modulo $p$.  Daher ist nur die \emph{halbe} Summe $S_p$ informativ.
\end{bemerkung}

\begin{theorem}[Wilson-Quotient via Halbfakultat]\label{satz:harmonisch}
  Für eine ungerade Primzahl $p$ gilt:
  \[
    W_p \;\equiv\; \mu_p + S_p \pmod{p}.
  \]
  Insbesondere ist $p$ genau dann eine Wilson-Primzahl, wenn
  $\mu_p + S_p \equiv 0 \pmod{p}$.
\end{theorem}

\begin{proof}
  Wir berechnen $(p-1)!$ modulo $p^2$.

  \textbf{Schritt 1 (Paarung).}
  Schreibe $(p-1)! = \prod_{j=1}^m j \cdot \prod_{j=m+1}^{p-1} j$.
  Substitution $j \mapsto p-j$ im zweiten Faktor:
  $\prod_{j=m+1}^{p-1} j = \prod_{j=1}^m (p-j)$.
  Also gilt:
  \[
    (p-1)! = \prod_{j=1}^m j(p-j).
  \]

  \textbf{Schritt 2 (Modulo $p^2$).}
  Für jedes $1 \leq j \leq m$ sei $j^{-1}$ das Inverse von $j$ modulo $p^2$.
  Dann gilt $p - j = -j(1 - pj^{-1})$ und somit $j(p-j) = -j^2(1-pj^{-1})$.
  Multipliziert über $j = 1, \ldots, m$:
  \[
    (p-1)! = (-1)^m (m!)^2 \cdot \prod_{j=1}^m (1 - pj^{-1}).
  \]
  Da $\prod_{j=1}^m (1-pj^{-1}) \equiv 1 - p\sum_{j=1}^m j^{-1} \pmod{p^2}$,
  ergibt sich:
  \[
    (p-1)! \;\equiv\; (-1)^m(m!)^2 \bigl(1 - pS_p\bigr) \pmod{p^2}.
  \]

  \textbf{Schritt 3 (Berechnung von $W_p$).}
  Nach dem Wilson-Satz gilt $(-1)^m(m!)^2 \equiv -1 \pmod{p}$,
  also schreibe $(-1)^m(m!)^2 = -1 + p\mu_p$ mit $\mu_p \in \Z$.
  Einsetzen:
  \begin{align*}
    (p-1)! &\equiv (-1+p\mu_p)(1-pS_p)\\
           &\equiv -1 + p\mu_p + pS_p - p^2\mu_pS_p\\
           &\equiv -1 + p(\mu_p + S_p) \pmod{p^2}.
  \end{align*}
  Daher $(p-1)! + 1 \equiv p(\mu_p + S_p) \pmod{p^2}$, und Division durch $p$:
  \[
    W_p \;=\; \frac{(p-1)!+1}{p} \;\equiv\; \mu_p + S_p \pmod p. \qedhere
  \]
\end{proof}

\begin{beispiel}
  Für $p = 5$: $m = 2$, $(m!)^2 = 4$, $(-1)^2 \cdot 4 = 4 = -1 + 5$,
  also $\mu_5 = 1$.  $S_5 = 1^{-1} + 2^{-1} = 1 + 3 = 4 \pmod 5$.
  $W_5 \equiv 1 + 4 = 5 \equiv 0 \pmod 5$. \checkmark\ ($5$ ist Wilson-Primzahl.)

  Für $p = 7$: $m = 3$, $(m!)^2 = 36$, $(-1)^3 \cdot 36 = -36 = -1 - 35 = -1 + 7(-5)$,
  also $\mu_7 = -5 \equiv 2 \pmod 7$.
  $S_7 = 1^{-1}+2^{-1}+3^{-1} = 1+4+5 = 10 \equiv 3 \pmod 7$.
  $W_7 \equiv 2+3 = 5 \pmod 7$. \checkmark\ ($W_7 = 103 \equiv 5 \pmod 7$, keine Wilson-Primzahl.)
\end{beispiel}

\begin{korollar}\label{kor:wilson_primzahl_kriterium_2}
  Eine ungerade Primzahl $p$ ist genau dann eine Wilson-Primzahl, wenn
  $\mu_p + S_p \equiv 0 \pmod{p}$.
\end{korollar}

\begin{proof}
  Direkte Folgerung aus Satz~\ref{satz:harmonisch}.
\end{proof}

% ================================================================
\section{Zusammenhang mit Bernoulli-Zahlen}
% ================================================================

\begin{theorem}[Wilson-Quotient und Bernoulli-Zahlen]\label{satz:bernoulli}
  Für eine ungerade Primzahl $p$ gilt:
  \[
    W_p \;\equiv\; -B_{p-1} \pmod{p},
  \]
  wobei $B_n$ die $n$-te Bernoulli-Zahl bezeichnet.
\end{theorem}

\begin{proof}[Beweis (Skizze über klassische Ergebnisse)]
  Die Aussage $W_p \equiv -B_{p-1} \pmod{p}$ ist ein klassisches Ergebnis,
  das den Wilson-Quotienten mit den $p$-adischen Bernoulli-Zahlen verbindet.
  Wir skizzieren die zwei wesentlichen Zutaten.

  \textbf{(A) Von Staudt--Clausen und $p$-adische Interpretation.}
  Der Satz von von Staudt--Clausen besagt:
  $B_{p-1} + \sum_{q\text{ prim},\,(q-1)\mid(p-1)} 1/q \in \Z$.
  Da $(p-1) \mid (p-1)$, tritt $p$ in dieser Summe auf, also hat $B_{p-1}$ den
  Nenner durch $p$ teilbar, und $pB_{p-1} \equiv -1 \pmod{p}$ im $p$-adischen Sinne.

  \textbf{(B) Zusammenhang mit der halben harmonischen Summe.}
  Aus Satz~\ref{satz:harmonisch} gilt $W_p \equiv \mu_p + S_p \pmod{p}$, wobei
  $S_p = \sum_{j=1}^{(p-1)/2} j^{-1}$ die \emph{halbe} harmonische Summe ist.
  (Die \emph{volle} Summe $\sum_{j=1}^{p-1} j^{-1} \equiv 0 \pmod{p}$ ist trivial
  und trägt keine Information über Wilson-Primzahlen.)
  Mithilfe der Potenzsummenformel $\sum_{j=1}^{p-1} j^k \equiv -B_k \pmod{p}$
  und Kummers Kongruenzen zeigt man, dass $\mu_p + S_p \equiv -B_{p-1} \pmod{p}$.
  Ein vollständiger Beweis findet sich in Ireland--Rosen \cite{Ireland1990}, Kap.~15.

  \textbf{Schlussfolgerung.}
  Eine Primzahl $p$ ist genau dann eine Wilson-Primzahl, wenn $p \mid W_p$,
  was nach (B) äquivalent zu $B_{p-1} \equiv 0 \pmod{p}$ ist.
\end{proof}

\begin{bemerkung}
  Der Satz von von Staudt--Clausen ist ein grundlegendes Ergebnis der
  Theorie der Bernoulli-Zahlen, das deren $p$-adische Eigenschaften präzise
  beschreibt. Die Kongruenz $W_p \equiv -B_{p-1} \pmod{p}$ verbindet die
  Wilson-Theorie mit der Theorie der unregelmäßigen Primzahlen von Kummer,
  die eine zentrale Rolle beim Beweis des letzten Satzes von Fermat spielte.
\end{bemerkung}

% ================================================================
\section{Bekannte Wilson-Primzahlen und Heuristik}
% ================================================================

\begin{theorem}[Bekannte Wilson-Primzahlen]\label{satz:bekannte_wilson}
  Die einzigen Wilson-Primzahlen unterhalb von $2 \times 10^{13}$ sind
  $p = 5$, $p = 13$ und $p = 563$.
\end{theorem}

\begin{proof}
  Verifikation für $p = 5$:
  $(5-1)! + 1 = 24 + 1 = 25 = 5^2$, also $5^2 \mid 25$. \checkmark

  Verifikation für $p = 13$:
  $(13-1)! + 1 = 479001600 + 1 = 479001601 = 13^2 \cdot 2834329$,
  also $13^2 \mid 479001601$. \checkmark

  Verifikation für $p = 563$:
  $W_{563} \equiv 0 \pmod{563}$ wurde durch direkte Computersuche verifiziert.

  Die Nichtexistenz weiterer Wilson-Primzahlen unterhalb von $2 \times 10^{13}$
  ist das Ergebnis einer erschöpfenden Computersuche von Crandall, Dilcher
  und Pomerance \cite{Crandall1997}.
\end{proof}

\begin{proposition}[Heuristische Dichte von Wilson-Primzahlen]\label{prop:heuristik}
  Heuristisch ist die Anzahl der Wilson-Primzahlen bis $x$ ungefähr $\ln \ln x$.
\end{proposition}

\begin{proof}[Beweis der Heuristik]
  Wir modellieren $W_p \bmod p$ als eine „zufällige" Restklasse in
  $\{0, 1, \ldots, p-1\}$. Die Wahrscheinlichkeit, dass eine gegebene
  Primzahl $p$ eine Wilson-Primzahl ist, beträgt dann ungefähr $1/p$.

  Nach dem Primzahlsatz ist die Anzahl der Primzahlen bis $x$ ungefähr
  $x/\ln x$. Die erwartete Anzahl von Wilson-Primzahlen bis $x$ ist daher
  \[
    \sum_{\substack{p \leq x \\ p \text{ prim}}} \frac{1}{p}
    \;\approx\; \ln \ln x,
  \]
  nach dem Mertenschen Satz $\sum_{p \leq x} 1/p = \ln \ln x + M + O(1/\ln x)$,
  wobei $M$ die Meissel--Mertens-Konstante ist.

  Diese Heuristik sagt voraus, dass Wilson-Primzahlen \emph{unendlich viele},
  aber \emph{sehr dünn gesät} sind: Bis $2 \times 10^{13}$ erwartet man
  ungefähr $\ln\ln(2 \times 10^{13}) \approx 3{,}4$ Wilson-Primzahlen,
  was mit den drei bekannten übereinstimmt.
\end{proof}

\begin{korollar}[Stärkere Kongruenz für Wilson-Primzahlen]\label{kor:staerker}
  Wenn $p$ eine Wilson-Primzahl ist, gilt $(p-1)! \equiv -1 \pmod{p^2}$.
\end{korollar}

\begin{proof}
  Nach Definition gilt $W_p = \bigl((p-1)!+1\bigr)/p$ und $p \mid W_p$.
  Das bedeutet $p^2 \mid (p-1)!+1$, d.\,h.\ $(p-1)! \equiv -1 \pmod{p^2}$.
\end{proof}

% ================================================================
\section{Schlussfolgerung}
% ================================================================

Der Wilson-Quotient $W_p$ ist ein reichhaltiges arithmetisches Invariant
von Primzahlen. Sein Zusammenhang mit der harmonischen Summe
$\sum j^{-1} \pmod{p}$ und mit der Bernoulli-Zahl $B_{p-1}$ ordnet ihn in
den Rahmen der $p$-adischen Analysis und Kummers Theorie der irregulären
Primzahlen ein.

Die drei bekannten Wilson-Primzahlen $5$, $13$, $563$ bleiben isolierte
Beispiele. Die heuristische Vorhersage von $\sim \ln\ln x$ Wilson-Primzahlen
bis $x$ legt ihre unendliche Existenz nahe, aber ein Beweis bleibt außer
Reichweite. Ob unendlich viele Wilson-Primzahlen existieren, ist eines der
klassischen offenen Probleme der analytischen Zahlentheorie.

\begin{thebibliography}{9}

\bibitem{Hardy1979}
G.~H.~Hardy und E.~M.~Wright.
\newblock \emph{Einführung in die Zahlentheorie} (An Introduction to the Theory of Numbers), 5.~Aufl.
\newblock Oxford University Press, 1979.

\bibitem{Ireland1990}
K.~Ireland und M.~Rosen.
\newblock \emph{A Classical Introduction to Modern Number Theory}, 2.~Aufl.
\newblock Springer, New York, 1990.

\bibitem{Crandall1997}
R.~Crandall, K.~Dilcher und C.~Pomerance.
\newblock A search for Wieferich and Wilson primes.
\newblock \emph{Math.\ Comp.}, 66(217):433--449, 1997.

\bibitem{Washington1997}
L.~C.~Washington.
\newblock \emph{Introduction to Cyclotomic Fields}, 2.~Aufl.
\newblock Springer, New York, 1997.

\bibitem{Berndt1994}
B.~C.~Berndt.
\newblock \emph{Ramanujan's Notebooks, Part IV}.
\newblock Springer, New York, 1994.

\end{thebibliography}

\end{document}
